\subsection{Задача Гурса}
\begin{theorem}[существования и единственности]
	Решение $u(x, y) \in C^2(\Pi) \cap C^1(\overline{\Pi})$ задачи Гурса:
	\[
	Lu = u_{xy} + au_x + b_y + cu = f(x, y),\ (x, y) \in \Pi = (0, l) \times (0, h) \quad (1.f)
	\]
	\[
	\textbf{Условия на характеристиках(УХ): } u|_{x=0} = \phi(y),\ 0 \le y 
	\le h
	u|_{y=0} = \psi(x),\ 0 \le x \le l 
	\]
	\[
	\textbf{Условия согласования(УС): } \phi(0) = \psi(o) = u(0, 0),
	\] 
	\[\text{если } a(x, y) \in C^1(\overline{\Pi}),\ b(x, y), c(x, y), f(x, y) \in C^1(\overline{\Pi}), \phi(y) \in C^2([0, h]),\ \psi \in C^2([0, l]),
	\] существует и единственно.
\end{theorem}

Физический смысл: "сопло Лаваля".

\subsection{Смешанная задача(СЗ) для однородного уравнения колебаний полубесконечной струны(ОУКПБС), закрепленной на левом конце.}

\begin{align*}
	& u_{tt} = a^2 u_{xx}, \quad x > 0,\ t > 0 \tag{1.0} \\
	\textbf{НУ: } \quad & u|_{t=0}=u_0(x),\ u_t|_{t=0}=u_1(x),\ x \ge 0 \tag{2} \\
	\textbf{ОГУ1Т: } \quad & u|_{x=0}=0,\ t \ge 0 \tag{3.1.0} \\
	\textbf{УС: } \quad & u_0(0)=u_1(0)=0=u_0^{''}(0), \tag{4.1.0}
\end{align*}

$u(x, t) \in C^2((0, +\infty) \times (0, \infty)) \cap C^1([0, +\infty) \times [0, \infty))$

$u_0(x) \in C^2(0, +\infty),\quad u_1(x) \in C^1(0, +\infty)$

Вспоминим Задачу Коши(ЗК) для ОУКБС:
\begin{align*}
	& u_{tt} = a^2 u_{xx}, \quad x > 0,\ t > 0 \tag{1.0} \\
	\textbf{НУ: } \quad & u|_{t=0}=u_0(x),\ u_t|_{t=0}=u_1(x),\ x \in \mathbb{R} \tag{2}
\end{align*}
$u(x, t) \in C^2(t>0) \cap C^1(t \ge 0),\quad u_0(x) \in C^2(\mathbb{R}),\quad u_1 \in C^1(\mathbb{R})$

\begin{lemma}
	Если НУ в ЗК (1.0), (2) являются нечётными функциями относительно т. $x=x_0$, то соответствующее решение $u(x, t)$ ЗК (1.0), (2) равно нулю в т. $x=x_0$, т.е. $u(x_0, t) = 0,\quad t \ge 0$.
\end{lemma}

\begin{lemma}
	Если НУ в ЗК (1.0), (2) являются чётными функциями относительно т. $x=x_0$, то частная производная соответствующего решение $u(x, t)$ ЗК (1.0), (2) равна нулю в т. $x=x_0$, т.е. $u_x(x_0, t) = 0,\quad t \ge 0$.
\end{lemma}

\begin{note}
	функция $y = f(x)$, определенная не $(a, b)$ и симметричная относительно точки $x=x_0 \in (a, b)$ называется чётной(нечётной), если $f(2x_0 - x) = f(x)(f(2x_0 - x) = -f(x))$, где $f(x_0 - x^{'}) = \pm f(x_0 + x^{'})$(относительно т.$x_0$).
\end{note}

\begin{proof}
	(Доказательство первой леммы)
	Положим $x_0 = 0$. Тогда по теореме о существовании и единственности решения ЗК (1.0), (2) для ОУКБС в рассматриваемых классах функций $\exists!$ и представляется Ф.Д.0(формулой Даламбера):
	\[
		u(x, t) = \frac{1}{2} [u_0(x-at) + u_0(x+at)] + \frac{1}{2a}\int_{x-at}^{x+at}u_1(\xi)d\xi \Rightarrow
	\]
	\[
		u(0, t) = \frac{u_0(-at)+u_0(at)}{2} + \frac{1}{2a} \int_{-at}^{at} u_1(\xi)d\xi = 0,
	\]
	при $t \ge 0$, т.к. из условий леммы функции $u_0(x)$ и $u_1(x)$ --- нечётные. 
	
	Аналогично, доказывается вторая лемма с учётом того факта, что если $f(x)$ --- чётная функция, то её производная $f^{'}(x)$ --- нечётная.
\end{proof}

\subsection{Метод продолжений}

\textbf{Следствие 1} Для того, чтобы свести СЗ (1.0), (2), (3.1.0), (4.1.0) для ОУКПБС, закрепленной на левом конце, к ЗК для ОУКБС нужно начальные функции $u_0(x)$ и $u_1(x)$ нечетным образом продолжить влево за т. $x=0$.

\textbf{Следствие 2} Для того, чтобы свести СЗ для ОУКПБС, со свободным левым концом, к ЗК для ОУКБС нужно начальные функции $u_0(x)$ и $u_1(x)$ четным образом продолжить влево за т. $x=0$.

\begin{theorem}
	Решение $u(x, t)$ СЗ (1.0), (2), (3.1.0), (4.1.0) для ОУКПБС, закрепленной на левом конце, существует и единственно и представляется обобщенной Ф.Д.0(Ф.Д.0.1):
	
\end{theorem}